% Part: proof-theory
% Chapter: cut-elimination
% Section: midsequent

\documentclass[../../../include/open-logic-section]{subfiles}

\begin{document}

\olfileid{pt}{cut}{mid}

\olsection{Teorema Sekèn Tengah lan Teorema Herbrand}

Salah siji akibat wigati saka teorema eliminasi cut yaiku
teorema sekèn tengah. Teorema iki mratélakaké yèn ana
!!a{proof}~$\pi$ kanggo~$\Gamma \Sequent \Delta$
kanthi saben !!{formula} preneks, banjur ana
!!a{proof}~$\pi'$ kang ngemot sekèn
$S \ident \Pi \Sequent \Lambda$ (diarani
\emph{sekèn tengah}) saéngga kabèh inferènsi ing
sadhuwuré $S$ ing~$\pi'$ iku proposisional lan kabèh
inferènsi ing sangisoré~$S$ iku inferènsi kuantor.
(!!^a{formula}~$!A$ diarani \emph{preneks} (utawa
nganggo \emph{wujud preneks}) yèn wujudé
\[\mathsf{Q}_1x_1\dots\mathsf{Q}_nx_n\,!B(x_1,\dots,x_n)\]
ing ngendi saben $\mathsf{Q}_i$ iku $\lforall$ utawa~$\lexists$.)
Akibat wigati saka teorema sekèn tengah yaiku teorema
Herbrand. Wujud umumé rada angèl diterangaké, nanging
kanthi garis gedhé: kanggo saben !!{provable}
!!{sentence}~$!A$ ing logika orde kapisan, ana tautologi
proposisional~$!A'$ kang bisa digunakaké kanggo
!!{prove}~$!A$. Iki versi kanggo !!{sentence}s kanthi
wujud~$\lexists[x][!B(x)]$ nalika $!B$ tanpa kuantor:

\begin{thm}[Teorema Herbrand]
Anggepen $!B(x)$ tanpa kuantor lan
$\Proves \lexists[x][!B(x)]$. Banjur ana suku-suku~$t_1$,
\dots, $t_n$ saéngga
$\Proves !B(t_1) \lor \dots \lor !B(t_n)$.
\end{thm}

!!{formula}~$!B(t_1) \lor \dots \lor !B(t_n)$ diarani
\emph{disjungsi Herbrand} saka~$\lexists[x][!B(x)]$.
Amarga $!B$ tanpa kuantor, disjungsi iki uga tanpa
kuantor. Mula yèn disjungsi mau !!{provable}, iku
!!{provable} nganggo bukti tanpa cut, lan akibaté
nganggo !!a{proof} kang mung ngemot inferènsi
proposisional. Dadi, disjungsi iku tautologi.

Pérangan kang angèl saka Teorema Herbrand yaiku nuduhaké
yèn saben !!{provable} !!{formula}~$\lexists[x][!B(x)]$
nduwèni disjungsi Herbrand. Arah kosok balènané gampang:
yèn $\lexists[x][!B(x)]$ nduwèni disjungsi Herbrand,
formula iku bisa dibuktèkaké. Pancèn, disjungsi Herbrand
ngakibataké $\lexists[x][!B(x)]$. Tuladhané, iki
!!a{proof} ing~$\Log{G3c}$ kanggo $n=2$:
\[
\Axiom$!B(t_1) \fCenter \lexists[x][!B(x)], !B(t_1), !B(t_2)$
\Axiom$!B(t_2) \fCenter \lexists[x][!B(x)], !B(t_1), !B(t_2)$
\RightLabel{\LeftR{\lor}}
\BinaryInf$!B(t_1) \lor !B(t_2) \fCenter \lexists[x][!B(x)], !B(t_1), !B(t_2)$
\RightLabel{\RightR{\lexists}}
\UnaryInf$!B(t_1) \lor !B(t_2) \fCenter \lexists[x][!B(x)], !B(t_2)$
\RightLabel{\RightR{\lexists}}
\UnaryInf$!B(t_1) \lor !B(t_2) \fCenter \lexists[x][!B(x)]$
\DisplayProof
\]

Kanggo njupuk disjungsi Herbrand saka
!!a{proof}~$\pi$ kanggo~$\Sequent \lexists[x][!B(x)]$,
gatekna kahanan nalika kabèh inferènsi $\RightR{\lexists}$
ing !!{proof}~$\pi$ tanpa cut ana ing pérangan pungkasan:
\[
\AxiomC{}
\RightLabel{$\pi_1$}
\Deduce$ \fCenter \lexists[x][!B(x)], !B(t_1), \dots, !B(t_n)$
\RightLabel{\RightR{\lexists}}
\UnaryInf$ \fCenter \lexists[x][!B(x)], !B(t_2), \dots, !B(t_n)$
\RightLabel{$(n-1) \times \RightR{\lexists}$}
\Deduce$ \fCenter \lexists[x][!B(x)]$
\DisplayProof
\]
Yèn iki kabèh inferènsi $\RightR{\lexists}$ ing~$\pi$
kang ngaktifaké $\lexists[x][!B(x)]$, banjur ora ana
inferènsi kuantor liya ing~$\pi_1$. Sebab, $!B(x)$
ora ngemot kuantor liya lan sekèn pungkasan ora ngemot
!!{formula}s ing sisih kiwa. Tegesé,
$ \fCenter \lexists[x][!B(x)], !B(t_1), \dots, !B(t_n)$
iku sekèn tengah saka~$\pi$. Kajaba iku, amarga ora
ana inferènsi kuantor ing sadhuwuré sekèn tengah,
saben sekèn ing sadhuwuré ngemot
$\lexists[x][!B(x)]$ minangka formula kang ora aktif
lan ora utama. Mula formula mau bisa dibusak saka
!!{proof}~$\pi_1$, lan kita ngolèhké !!a{proof}
kanggo~$\fCenter !B(t_1), \dots, !B(t_n)$.
Kanthi $n-1$ panrapan $\RightR{\lor}$, kita banjur
ngolèhké !!a{proof} kanggo disjungsi Herbrand
~$\Sequent !B(t_1) \lor \dots \lor !B(t_n)$.

Masalahé, sanadyan !!{proof} tanpa cut kanggo
$\Sequent \lexists[x][!B(x)]$, kita ora bisa nganggep
kabèh inferènsi \RightR{\lexists} ana ing siji runtutan
ing pungkasan. Bisa ana inferènsi ing antarané inferènsi
\RightR{\lexists} mau kang sawetara $!B(t_i)$ dadi
formula utama. Mula teorema sekèn tengah dibutuhaké.

\begin{thm}[Teorema Sekèn Tengah]\ollabel{thm:midsequent}
Yèn ana \Log{G3c}-!!{proof}~$\pi$ kanggo
$\Gamma \Sequent \Delta$ kanthi kabèh !!{formula}s ing
$\Gamma$ lan $\Delta$ preneks, banjur ana
!!a{proof}~$\pi'$ kang ngemot sekèn
$\Pi \Sequent \Lambda$ saéngga kabèh inferènsi
ing sangisoré iku inferènsi kuantor lan kabèh
inferènsi ing sadhuwuré iku inferènsi proposisional.
\end{thm}

\begin{proof}
  Titepna \emph{urutan} inferènsi kuantor ing~$\pi$
  minangka cacah inferènsi proposisional ing sangisoré,
  lan urutan~$o(\pi)$ saka $\pi$ minangka jumlah
  urutan kabèh inferènsi kuantor ing kono.
  Yèn $o(\pi) = 0$, ora ana inferènsi kuantor kang
  nduwèni inferènsi proposisional ing sangisoré,
  mula asilé wis kacukupan. Amarga saben inferènsi
  kuantor nduwèni siji premis, yèn ana inferènsi
  kuantor, ana siji kang paling dhuwur lan premisé
  dadi sekèn tengah. Yèn ora ana inferènsi kuantor,
  sekèn pungkasané dhéwé dadi sekèn tengah.

  Yèn $o(\pi) > 0$, pilihen inferènsi kuantor paling
  ngisor kang urutané $>0$. Amarga urutané $>0$,
  mesthi ana inferènsi proposisional ing sangisoré.
  Amarga inferènsi kuantor iki paling ngisor,
  konklusiné ora bisa dadi premis inferènsi kuantor:
  yèn iya, inferènsi kuantor kapindho kang luwih
  ngisor iku uga nduwèni inferènsi proposisional ing
  sangisoré. Kita bedakaké (akèh!) kahanan manut
  inferènsi kuantor lan inferènsi proposisional ing
  sangisoré. Amarga sekèn pungkasan mung ngemot
  formula preneks, formula utama inferènsi kuantor
  ora bisa aktif ing inferènsi proposisional sabanjuré.
  Yèn aktif, formula utama inferènsi proposisional iku
  bakal nduwèni kuantor ing sangisoré konektif
  proposisional. Miturut sipat sub-!!{formula},
  iki kudu dadi sub-!!{formula} saka !!a{formula}
  ing sekèn pungkasan. Mula !!{formula} utama
  inferènsi kuantor dadi !!a{element} kontèks inferènsi
  proposisional kang luwih ngisor.

  Iki rong tuladha:

Sepisan, anggepen inferènsi $\RightR{\lforall}$ ana ing
premis tengen saka $\LeftR{\lif}$. Banjur
sub-!!{proof} kang cocog ing~$\pi$ pungkasané ana ing
sub-!!{proof}~$\pi_1$.
\[
\AxiomC{}
\RightLabel{$\pi_2$}
\Deduce$\Gamma' \fCenter \Delta', \lforall[x][!A(x)], !B$
\AxiomC{}
\RightLabel{$\pi_3$}
\Deduce$!C, \Gamma' \fCenter \Delta', !A(c)$
\RightLabel{\RightR{\lforall}}
\UnaryInf$!C, \Gamma' \fCenter \Delta', \lforall[x][!A(x)]$
\RightLabel{\LeftR{\lif}}
\BinaryInf$!B \lif !C, \Gamma' \fCenter \Delta', \lforall[x][!A(x)]$
\DisplayProof
\]
Kita bisa nganggep $\pi$ reguler, mula variabel
eigen~$c$ ora ana ing sanjabané~$\pi_3$;
mliginé, ora ana ing $!B \lif !C$ utawa ing~$\pi_2$.
Saiki gatekna bukti~$\pi_1'$:
\[
\AxiomC{}
\RightLabel{$\pi_2'$}
\Deduce$\Gamma' \fCenter \Delta', \lforall[x][!A(x)], !A(c), !B$
\AxiomC{}
\RightLabel{$\pi_3$}
\Deduce$!C, \Gamma' \fCenter \Delta', \lforall[x][!A(x)], !A(c)$
\RightLabel{\LeftR{\lif}}
\BinaryInf$!B \lif !C, \Gamma' \fCenter \Delta', \lforall[x][!A(x)], !A(c)$
\RightLabel{\RightR{\lforall}}
\UnaryInf$!B \lif !C, \Gamma' \fCenter \Delta', \lforall[x][!A(x)], \lforall[x][!A(x)]$
\DisplayProof
\]
Ing kéné $\pi_2'$ diolehké saka~$\pi_2$ kanthi
pelemahan nambah~$!A(c)$ ing tengen. (Iki ora nambah
inferènsi apa waé, mliginé ora nambah inferènsi kuantor
kang bisa ngowahi urutan. Iki uga ora nerak syarat
variabel eigen ing~$\pi_2$, amarga ora ana konstanta
ing~$!A(c)$ kang dadi variabel eigen ing~$\pi_2$.)
Syarat variabel eigen kanggo \RightR{\lforall} tetep
kacukupan, amarga $c$ ora ana ing~$!B \lif !C$.

Kita migunakaké admisibilitas kontraksi kanggo
ngolèhké !!a{proof}~$\pi_1''$ kanggo
$!B \lif !C, \Gamma' \fCenter \Delta', \lforall[x][!A(x)]$.
Bukti admisibilitas $\RightR{\Contraction}$ kanggo
$\lforall[x][!A(x)]$ lan bukti inversibilitas
$\RightR{\lforall}$ kang digunakaké ing kono ora
ngowahi urutan inferènsi apa waé: paling-paling
mung mbusak inferènsi. Mula cacah inferènsi kuantor
ing~$\pi_1''$ ora luwih akèh tinimbang ing~$\pi_1$,
lan saben inferènsi kuantor ing~$\pi_1''$ nduwèni
cacah inferènsi proposisional ing sangisoré kang padha
karo inferènsi kang cocog ing~$\pi_1$---kajaba
\RightR{\lforall} pungkasan: ana
\LeftR{\lif} ing sangisoré ing~$\pi$, déné ing~$\pi_1''$
wis dipindhah ing sadhuwuré. Gantenana $\pi_1$ ing
$\pi$ nganggo $\pi_1''$. Urutan !!{proof} asilé
luwih cilik, amarga paling ora urutan inferènsi
$\RightR{\lforall}$ wis suda (paling ora)~$1$.
Saiki trapna hipotesis induksi.

Kahanan nalika inferènsi kuantor iku
\RightR{\lexists} malah luwih gampang, amarga kontraksi
ora dibutuhaké lan syarat variabel eigen ora perlu
digatèkké. Tuladhané, anggepen $\RightR{\lexists}$
disusul $\RightR{\land}$ (saiki minangka premis kiwa).
Sub-!!{proof} kang cocog yaiku $\pi_1$:
\[
\AxiomC{}
\RightLabel{$\pi_2$}
\Deduce$\Gamma' \fCenter \Delta', \lexists[x][!A(x)], !A(t), !B$
\RightLabel{\RightR{\lexists}}
\UnaryInf$\Gamma' \fCenter \Delta', \lexists[x][!A(x)], !B$
\AxiomC{}
\RightLabel{$\pi_3$}
\Deduce$\Gamma' \fCenter \Delta', \lexists[x][!A(x)], !C$
\RightLabel{\RightR{\land}}
\BinaryInf$\Gamma' \fCenter \Delta', \lexists[x][!A(x)], !B \land !C$
\DisplayProof
\]
Gantenana $\pi_1$ ing~$\pi$ nganggo $\pi_1'$:
\[
\AxiomC{}
\RightLabel{$\pi_2$}
\Deduce$\Gamma' \fCenter \Delta', \lexists[x][!A(x)], !A(t), !B$
\AxiomC{}
\RightLabel{$\pi_3'$}
\Deduce$\Gamma' \fCenter \Delta', \lexists[x][!A(x)], !A(t), !C$
\RightLabel{\RightR{\land}}
\BinaryInf$\Gamma' \fCenter \Delta', \lexists[x][!A(x)], !A(t), !B \land !C$
\RightLabel{\RightR{\lexists}}
\UnaryInf$\Gamma' \fCenter \Delta', \lexists[x][!A(x)], !B \land !C$
\DisplayProof
\]
Ing kéné $\pi_3'$ diolehké saka $\pi_3$ kanthi
pelemahan nambah~$!A(t)$ ing tengen. Pelemahan iki
mung nambah $!A(t)$ ing sisih tengen saben sekèn
ing~$\pi_3$, mula ora ngowahi urutan inferènsi kuantor.
Urutan inferènsi kuantor ing~$\pi'$ padha karo
ing~$\pi$, kajaba urutan inferènsi \RightR{\lexists}
kang diowahi urutané ngliwati~\RightR{\land}
suda~$1$. Hipotesis induksi bisa ditrapaké.
\end{proof}

\begin{prob}
Jlentrehna kahanan ing bukti \olref{thm:midsequent}
nalika \LeftR{\lexists} ana ing premis~\RightR{\lif},
lan nalika \LeftR{\lforall} ana ing premis kiwa~\LeftR{\lor}.
\end{prob}

\begin{proof}[Bukti Teorema Herbrand]
Anggepen $\pi$ pungkasané ana ing
$\Sequent \lexists[x][!A(x)]$. Miturut
\olref{thm:midsequent}, $\pi$ bisa diowahi dadi
!!a{proof}~$\pi'$ kang nduwèni sekèn tengah~$S$,
kang kudu wujudé \[\Sequent
\lexists[x][!A(x)], !A(t_1), \dots, !A(t_n).\]
Amarga kabèh inferènsi ing sadhuwuré $S$ proposisional,
$\lexists[x][!A(x)]$ ora bisa dadi formula utama
inferènsi ing sadhuwuré~$S$. Amarga ora atomik,
formula iki uga ora bisa dadi formula utama ing aksioma.
Amarga $\lexists[x][!A(x)]$ ora bisa dadi subformula
sejati saka $!A(t_1)$ (élinga $!A(x)$ tanpa kuantor),
formula iki uga ora bisa aktif ing sadhuwuré~$S$.
Mula mbusak $\lexists[x][!A(x)]$ saka
sub-!!{proof} kang pungkasané ana ing~$S$
ngasilaké !!{proof} kang bener kanggo
$\Sequent !A(t_1), \dots, !A(t_n)$; saka kéné,
kita ngolèhké !!a{proof} kanggo
$\Sequent !A(t_1) \lor \dots \lor !A(t_n)$
kanthi $n-1$ inferènsi $\RightR{\lor}$.
\end{proof}

% \begin{prob}
% Find a Herbrand disjunction of $\lexists[x][\lforall[y][(!A(x) \lif
% !A(y))]]$ by finding a cut-free !!{proof} of $\Sequent
% \lexists[x][\lforall[y][(!A(x) \lif !A(y))]]$ and applying the
% transformation in \olref{thm:midsequent} to find a midsequent (if
% necessary).
% \end{prob}

\end{document}
